% =====================================================================
%  APENDICES A -- E
% =====================================================================

\newpage
\appendix
\section{Plan del proyecto de Ingeniería de Requisitos}

\subsection{Identificación de las partes interesadas}

\begin{longtable}{|C{1.6cm}|L{4.4cm}|C{1.8cm}|C{1.6cm}|C{1.6cm}|C{2.6cm}|}
\hline
\rowcolor{grisclaro}\textbf{Código} & \textbf{Rol} & \textbf{Organización} & \textbf{Tipo} & \textbf{Influencia} & \textbf{Evidencia} \\
\hline
\endhead
\id{ENTR-01} & Administrador general & Palmicultora M & Primario & Alta & \id{EV-01} \\ \hline
\id{ENTR-02} & Asesor técnico & Palmicultora M & Primario & Alta & \id{EV-02} \\ \hline
\id{ENTR-03} & Jefe de polinización & Palmicultora M & Primario & Alta & \id{EV-03} \\ \hline
\id{ENTR-04} & Supervisor de técnicos agrícolas & Extractora R & Externo & Alta & \id{EV-04} \\ \hline
\id{ENTR-05} & Trabajador agrícola & Palmicultora M & Operativo & Baja & \id{EV-05} \\ \hline
\id{ENTR-06} & Trabajador agrícola & Palmicultora M & Operativo & Baja & \id{EV-06} \\ \hline
\id{ENTR-07} & Asistente de administración & Palmicultora M & Primario & Media & \id{EV-07} \\ \hline
\id{ENTR-08} & Técnico agrícola & Extractora R & Externo & Media & \id{EV-08} \\ \hline
\caption{Identificación de las partes interesadas entrevistadas}
\end{longtable}

\subsection{Cronograma de actividades}

\begin{longtable}{|C{2.0cm}|L{6.6cm}|L{6.4cm}|}
\hline
\rowcolor{grisclaro}\textbf{Semana} & \textbf{Actividad} & \textbf{Producto} \\
\hline
\endhead
1--3 & Identificación del problema y del sistema objetivo & Propuesta de proyecto \\ \hline
3--4 & Primera ronda de elicitación (\id{EV-01} a \id{EV-03}) y observación de campo & Entrega 1 (1A) \\ \hline
5--6 & Elaboración del paquete ético y obtención del aval institucional & Documentos A.1--A.13 y Categoría C \\ \hline
6--9 & Modelado UML inicial y formalización de requisitos & Diagramas y plantillas de RF \\ \hline
9--10 & Consolidación del ERS parcial & Entrega 2 (1B) \\ \hline
11 & Adenda ética y registro del protocolo experimental en OSF & Adenda y comprobante OSF \\ \hline
12 & Segunda ronda de campo (\id{EV-04} a \id{EV-08}) y ampliación del cuestionario & Grabaciones, transcripciones, $n=62$ \\ \hline
13 & Consolidación del ERS completo y del MVP & \textbf{Entrega 3 (2A)} \\ \hline
14--16 & Ejecución del componente empírico y análisis & Resultados y borrador de manuscrito \\ \hline
17 & Defensa final & Entrega 4 (2B) \\ \hline
\caption{Cronograma de actividades de ingeniería de requisitos}
\end{longtable}

\subsection{Técnicas de elicitación y justificación}

\begin{longtable}{|L{3.6cm}|C{1.6cm}|L{9.4cm}|}
\hline
\rowcolor{grisclaro}\textbf{Técnica} & \textbf{Sesiones} & \textbf{Justificación} \\
\hline
\endhead
Entrevista semiestructurada & 8 & Permite profundizar en el conocimiento experto y adaptarse a respuestas inesperadas. Es la técnica indicada cuando el dominio contiene conocimiento tácito \cite{ferrari2016}, situación confirmada en el criterio de madurez del racimo. \\ \hline
Cuestionario estructurado & 2 aplicaciones & Permite cuantificar la utilidad percibida y las condiciones de operación en una población que no puede ser entrevistada en su totalidad. Instrumento idéntico entre rondas para preservar la comparabilidad. \\ \hline
Observación directa & 1 & Documenta el proceso tal como ocurre, sin la mediación del relato. Registrada en \id{EV-09}. \\ \hline
Análisis de documentos & 2 tipos & Los planes semanales (\id{EV-11}) revelaron la estructura de planificación y las unidades reales de medida del trabajo. La hoja de liquidación (\id{EV-13}) hizo explícito el catálogo de tarifas por labor, que ninguno de los ocho participantes enunció pese a operar bajo él. Es la técnica que más requisitos aportó por unidad de esfuerzo. \\ \hline
Prototipado & 8 pantallas & Los prototipos de alta fidelidad (\id{MU-01} a \id{MU-08}) permiten validar la interfaz con personal de baja alfabetización tecnológica antes de construir. \\ \hline
\caption{Técnicas de elicitación aplicadas y su justificación}
\end{longtable}

\subsection{Declaración de conflicto de interés}

Se declara que dos de las personas participantes mantienen relación de consanguinidad con el
analista líder del equipo: el participante \id{ENTR-02} en primer grado y el participante
\id{ENTR-07} en segundo grado. Ambas entrevistas fueron conducidas por el propio analista líder.

Se aplicaron seis medidas de mitigación:

\begin{enumerate}[nosep]
  \item Presencia de un segundo integrante del equipo en calidad de registrador durante ambas sesiones.
  \item Conservación íntegra de las grabaciones y transcripciones, disponibles para auditoría docente.
  \item Uso estricto del banco de preguntas aprobado, sin preguntas improvisadas fuera del guion.
  \item Verificación de los requisitos derivados por integrantes con rol de verificador que no participaron en esas sesiones.
  \item Triangulación obligatoria de todo hallazgo con al menos una de las otras seis entrevistas.
  \item \textbf{Ningún requisito del presente ERS se sustenta de forma exclusiva en estas dos fuentes.}
\end{enumerate}

\noindent
La sexta medida es verificable en la matriz de trazabilidad: toda fila que referencia \id{EV-02} o
\id{EV-07} referencia además al menos una evidencia adicional independiente.

% =====================================================================
\newpage
\section{Trabajo de campo y evidencias}

\subsection{Catálogo de evidencias}

\footnotesize
\begin{longtable}{|C{1.2cm}|L{2.6cm}|C{1.7cm}|C{1.8cm}|L{4.2cm}|C{1.3cm}|C{1.1cm}|}
\hline
\rowcolor{grisclaro}\textbf{ID-EV} & \textbf{Tipo} & \textbf{Fecha} & \textbf{Participante} & \textbf{Requisitos derivados} & \textbf{Zona} & \textbf{Cons.} \\
\hline
\endhead
\id{EV-01} & Entrevista (audio) & 23/05/2026 & \id{ENTR-01} & \id{RF-01} a \id{RF-06}, \id{RF-09}, \id{RF-12}, \id{RF-15} a \id{RF-20} & [R] & Sí \\ \hline
\id{EV-02} & Entrevista (video) & 23/05/2026 & \id{ENTR-02} & \id{RF-02}, \id{RF-05}, \id{RF-08}, \id{RF-11}, \id{RF-17}, \id{RF-18} & [R] & Sí \\ \hline
\id{EV-03} & Entrevista (video) & 23/05/2026 & \id{ENTR-03} & \id{RF-13}, \id{RF-14}, \id{RF-15}, \id{RNF-05} & [R] & Sí \\ \hline
\id{EV-04} & Entrevista (video) & 28/07/2026 & \id{ENTR-04} & \id{RF-10}, \id{RF-21} a \id{RF-25}, \id{RF-32}, \id{RF-33}, \id{RNF-18} & [R] & Sí \\ \hline
\id{EV-05} & Entrevista (video) & 28/07/2026 & \id{ENTR-05} & \id{RF-04}, \id{RF-28}, \id{RF-35} & [R] & Sí \\ \hline
\id{EV-06} & Entrevista (video) & 28/07/2026 & \id{ENTR-06} & \id{RF-04}, \id{RF-28}, \id{RF-29} & [R] & Sí \\ \hline
\id{EV-07} & Entrevista (video) & 02/08/2026 & \id{ENTR-07} & \id{RF-29} a \id{RF-31}, \id{RNF-11} & [R] & Sí \\ \hline
\id{EV-08} & Entrevista (video) & 28/07/2026 & \id{ENTR-08} & \id{RF-21}, \id{RF-22}, \id{RF-24}, \id{RD-07} & [R] & Sí \\ \hline
\id{EV-09} & Observación de campo & 23/05/2026 & Equipo AHMRV & \id{RF-01}, \id{RF-02}, \id{RF-05} & [P] & N/A \\ \hline
\id{EV-10} & Cuestionario piloto & 30/06/2026 & $n=4$ & Validación de \id{RF-04}, \id{RF-07}, \id{RF-12} & [P] & Anónimo \\ \hline
\id{EV-11} & Documentos de la organización & 29/06 a 24/07/2026 & Palmicultora M & \id{RF-02}, \id{RF-03}, \id{RF-19}, \id{RF-26} a \id{RF-28}, \id{RF-34} & [P] + [R] & N/A \\ \hline
\id{EV-12} & Cuestionario ampliado & 30/06/2026 & $n=62$ & \id{RF-12}, \id{RF-15}, \id{RF-35}, \id{RD-10}, \id{RNF-08} & [P] & Anónimo \\ \hline
\id{EV-13} & Hoja de liquidación semanal & 03--07/08/2026 & Palmicultora M & \id{RF-36} a \id{RF-39} & [P] + [R] & N/A \\ \hline
\caption{Catálogo de evidencias de campo. [P] = zona pública, [R] = zona restringida cifrada}
\label{tab:catalogo}
\end{longtable}

\noindent
\textbf{Nota sobre la renumeración.} Las evidencias \id{EV-09} y \id{EV-10} correspondían a
\id{EV-04} y \id{EV-05} en la Entrega 2 (1B). El desplazamiento permitió asignar identificadores
contiguos a las entrevistas de la segunda ronda. La equivalencia se registra en \id{CHANGELOG.md}.

\subsection{Documentos de la organización}

Se obtuvieron dos tipos documentales distintos, ambos correspondientes a la operación real de la
organización cliente durante el período del proyecto.

\begin{longtable}{|C{1.4cm}|L{4.2cm}|C{2.4cm}|L{6.4cm}|}
\hline
\rowcolor{grisclaro}\textbf{ID-EV} & \textbf{Tipo documental} & \textbf{Período} & \textbf{Aporte a la especificación} \\
\hline
\endhead
\id{EV-11} & Plan semanal de labores (presupuesto por cumplir y cumplido) & Semanas 27 a 30 de 2026 & Estructura de lotes, catálogo de labores, unidades de avance, patrón de arrastre de labor incompleta, requerimientos de insumos \\ \hline
\id{EV-13} & Hoja de liquidación de presupuesto contra ejecución & Semana 32 de 2026 & Código de labor, unidad de medida, tarifa vigente, distinción entre cantidad facturada y ejecutada, grupo de labor no presupuestada, flete como costo independiente, aprobación de la administración \\ \hline
\caption{Documentos de la organización incorporados como evidencia}
\end{longtable}

\noindent
\textbf{Relación entre ambos documentos.} \id{EV-11} declara metas; \id{EV-13} liquida su ejecución.
El par constituye el ciclo completo de planificación y control que el sistema debe reproducir, y
sustenta los requisitos \id{RF-26}, \id{RF-27} y \id{RF-37} a \id{RF-39}. Su obtención en momentos
distintos del proyecto explica que los requisitos de liquidación aparezcan al final de la numeración
pese a corresponder a un proceso central de la organización.

\noindent
\textbf{Tratamiento aplicado.} Conforme a la §8.4 de la guía, la copia depositada en la zona pública
tiene tachados de forma irreversible los valores de tarifa, los importes por labor y los totales de
gasto, así como el nombre propio que figuraba en la columna de observaciones, sustituido por el
código de participante correspondiente. Los originales íntegros residen en el contenedor cifrado.

\subsection{Instrumentos aplicados}

Los instrumentos residen en \id{06\_Experimento/instrumentos/}. Se aplicaron cuatro guiones
diferenciados por perfil, decisión metodológica que permite ajustar el vocabulario y la profundidad
técnica al interlocutor:

\begin{longtable}{|C{2.0cm}|L{5.0cm}|C{2.0cm}|L{5.4cm}|}
\hline
\rowcolor{grisclaro}\textbf{Guion} & \textbf{Perfil destinatario} & \textbf{Aplicado a} & \textbf{Foco temático} \\
\hline
\endhead
\id{G-01} & Dirección y asesoría técnica & \id{EV-01}, \id{EV-02} & Manejo del cultivo, nutrición, sanidad, estimación \\ \hline
\id{G-02} & Jefatura de polinización & \id{EV-03} & Antesis, calidad de la labor, control del personal \\ \hline
\id{G-03} & Personal de la extractora & \id{EV-04}, \id{EV-08} & Criterios de calidad, calificación, penalización \\ \hline
\id{G-04} & Personal de campo y asistencia & \id{EV-05} a \id{EV-07} & Jornada, registro, reporte de afectaciones, tecnología \\ \hline
\id{C-01} & Cuestionario (idéntico en ambas rondas) & \id{EV-10}, \id{EV-12} & Perfil, registro actual, dificultades, utilidad percibida, aceptación \\ \hline
\caption{Instrumentos de recolección aplicados}
\end{longtable}

\subsection{Codificación temática y saturación}

La codificación temática de las ocho transcripciones reside en
\id{07\_Datos/datos\_crudos/}\allowbreak\id{codificacion.csv}, con las columnas Fragmento, Código,
Categoría, Requisito derivado, ID de evidencia y Analista codificador.

\begin{longtable}{|C{1.6cm}|C{2.4cm}|C{2.4cm}|C{2.4cm}|L{5.0cm}|}
\hline
\rowcolor{grisclaro}\textbf{Entrevista} & \textbf{Códigos totales} & \textbf{Códigos nuevos} & \textbf{Acumulado} & \textbf{Observación} \\
\hline
\endhead
\id{EV-01} & 24 & 24 & 24 & Primera fuente: todo es nuevo \\ \hline
\id{EV-02} & 21 & 10 & 34 & Alta redundancia con \id{EV-01}: 52\% de los fragmentos reiteran códigos previos \\ \hline
\id{EV-03} & 14 & 6 & 40 & Aporta el dominio de la polinización asistida \\ \hline
\id{EV-04} & 26 & \textbf{17} & 57 & \textbf{Repunte}: dominio nuevo (calidad de la fruta en la extractora) \\ \hline
\id{EV-05} & 9 & 4 & 61 & Confirma el registro manual en papel \\ \hline
\id{EV-06} & 11 & 1 & 62 & \textbf{Mínimo de la serie}: 9\% de aportación \\ \hline
\id{EV-07} & 15 & 4 & 66 & Aporta la latencia de verificación y el canal informal \\ \hline
\id{EV-08} & 18 & 2 & 68 & Triangula \id{EV-04}; 11\% de aportación \\ \hline
\caption{Aportación de códigos nuevos por entrevista (curva de saturación)}
\label{tab:saturacion}
\end{longtable}

\begin{figure}[htbp]
\centering
\includegraphics[width=0.96\textwidth]{img/curva_saturacion}
\caption{Curva de saturación temática: aportación de códigos nuevos por entrevista, en orden cronológico de recolección.}
\label{fig:saturacion}
\end{figure}

\noindent
\textbf{Lectura de la curva.} La aportación de códigos nuevos desciende de forma sostenida hasta
\id{EV-03}, con \id{EV-02} reiterando el 52\% de sus fragmentos y \id{EV-03} el 57\%. Ese descenso
sugería saturación al cierre de la primera ronda.

El repunte de \id{EV-04} (17 códigos nuevos sobre 26 fragmentos, un 65\% de aportación) demuestra
que dicha saturación era \textbf{aparente}: correspondía al agotamiento de un dominio ---el manejo
del cultivo--- y no del problema completo. La incorporación de la organización externa abrió un
dominio no explorado: los criterios de calidad con que se valora económicamente la producción.

La caída posterior es igual de informativa. \id{EV-06} aporta un solo código nuevo sobre once
fragmentos (9\%) y \id{EV-08} aporta dos sobre dieciocho (11\%), pese a pertenecer este último al
mismo dominio recién abierto. Es decir, el dominio de la calidad de la fruta \emph{sí} alcanzó
saturación, pero con dos fuentes, no con ocho entrevistas en total.

Este comportamiento ---saturación por dominio y no por número absoluto de entrevistas--- es
coherente con la distinción entre saturación de código y saturación de significado de Hennink
\emph{et al.} \cite{hennink2017}. Se reportará como hallazgo metodológico en el manuscrito, con la
implicación práctica de que el criterio de suficiencia muestral debe evaluarse por dominio de
conocimiento y no sobre el agregado.

\subsection{Zonas de evidencia y verificación}

Conforme a la §4.1 de la guía, la evidencia se organiza en dos zonas. La zona pública \textbf{[P]}
contiene transcripciones anonimizadas, copias enmascaradas de consentimientos, fotografías sin
rostros ni coordenadas GPS, respuestas del cuestionario sin columnas identificativas y las fichas
técnicas de cada archivo multimedia. La zona restringida \textbf{[R]} contiene los originales
identificables, cifrados con AES-256 en
\id{02\_Evidencias/00\_Restringido/}\allowbreak\id{evidencias\_restringidas\_2A.7z}.

El hash SHA-256 de cada archivo se calculó \emph{antes} del cifrado y consta en
\id{checksums.sha256} junto con la ficha técnica pública \id{fichas\_tecnicas.csv}, que registra
nombre de archivo, tipo, fecha, código de participante, duración, códec, tamaño y hash. La
contraseña del contenedor se entrega únicamente al docente por el espacio de la actividad en el
sistema de gestión académica.

\noindent\fbox{\parbox{0.97\textwidth}{%
\small\textbf{Nota técnica sobre Git LFS.} Los archivos multimedia y el contenedor cifrado se
versionan mediante Git Large File Storage. Una clonación sin \id{git lfs pull} muestra archivos de
133 bytes correspondientes a los punteros, no al contenido. Las instrucciones de recuperación
constan en el \id{README.md} de la raíz del repositorio.}}

% =====================================================================
\newpage
\section{Clasificación y priorización completa}

La tabla completa de clasificación de los 42 requisitos funcionales, con su prioridad MoSCoW, su
categoría de Kano, los cuatro factores del cálculo WSJF y la justificación de cada asignación,
reside en \id{04\_Trazabilidad/}\allowbreak\id{priorizacion\_moscow\_kano.csv}. La §5.3 del cuerpo del documento
desarrolla la argumentación de las categorías extremas y el contraste entre las tres técnicas.

Los requisitos no funcionales y las restricciones de diseño se clasifican como sigue.

\begin{longtable}{|C{2.0cm}|C{2.0cm}|L{5.0cm}|L{5.4cm}|}
\hline
\rowcolor{grisclaro}\textbf{ID} & \textbf{Prioridad} & \textbf{Característica / Ámbito} & \textbf{Razón} \\
\hline
\endhead
\id{RNF-01}, \id{RNF-02} & Must & Adecuación funcional & La exactitud del clasificador determina la utilidad del sistema \\ \hline
\id{RNF-03} & Should & Eficiencia de desempeño & Deseable; sin umbral crítico en la operación de campo \\ \hline
\id{RNF-04} & Must & Eficiencia de desempeño (IA) & Un diagnóstico lento no se usa en campo \\ \hline
\id{RNF-05} & Should & Eficiencia (precisión GPS) & Condiciona la confiabilidad del conteo \\ \hline
\id{RNF-06} & Should & Compatibilidad & Requerido para el intercambio con la extractora \\ \hline
\id{RNF-07} & Should & Compatibilidad & Cobertura de versiones de plataforma \\ \hline
\id{RNF-08}, \id{RNF-09} & Must & Interacción de usuario & Condicionan la adopción por personal de baja alfabetización digital \\ \hline
\id{RNF-10} & Should & Fiabilidad & Disponibilidad alta; no bloqueante en v1 \\ \hline
\id{RNF-11} & Must & Fiabilidad & El 74{,}2\% del personal carece de conectividad permanente \\ \hline
\id{RNF-12} & Should & Fiabilidad & Tolerancia a degradación de señal GPS \\ \hline
\id{RNF-13}, \id{RNF-14} & Must & Seguridad & Exigidos por los artículos 37 y 40 de la LOPDP \\ \hline
\id{RNF-19} & Must & Seguridad física & Condiciona la recomendación fitosanitaria de \id{RF-09}: su omisión expone físicamente al personal \\ \hline
\id{RNF-15} & Could & Mantenibilidad & Meta de calidad interna \\ \hline
\id{RNF-16} & Must & Explicabilidad & Obligatorio para todo componente de IA; gatekeeper G7 \\ \hline
\id{RNF-17} & Should & Portabilidad & Evita dependencia de un proveedor de inferencia \\ \hline
\id{RNF-18} & Must & Flexibilidad & Los umbrales varían por variedad; prerrequisito de \id{RF-21} \\ \hline
\id{RD-01} a \id{RD-10} & Must & Restricciones & Impuestas por el entorno, la normativa o la contraparte \\ \hline
\caption{Clasificación de los requisitos no funcionales y las restricciones}
\end{longtable}

% =====================================================================
\newpage
\section{Matriz de trazabilidad completa}

La matriz completa, con las trece columnas exigidas por la §8.8 de la guía, reside en
\id{04\_Trazabilidad/}\allowbreak\id{matriz\_trazabilidad.csv} en formato procesable. La §5.4 del cuerpo del
documento presenta la vista reducida de sus 55 filas y el diagnóstico de cobertura correspondiente.

\subsection{Diccionario de identificadores}

\begin{longtable}{|C{2.2cm}|L{5.2cm}|L{7.0cm}|}
\hline
\rowcolor{grisclaro}\textbf{Prefijo} & \textbf{Significado} & \textbf{Rango} \\
\hline
\endhead
\id{TR-XX} & Fila de la matriz de trazabilidad & \id{TR-01} a \id{TR-55} \\ \hline
\id{CB-XX} & Criterio de aceptación BDD derivado de la ficha del requisito & \id{CB-01} a \id{CB-37} \\ \hline
\id{O-X} & Objetivo del sistema & \id{O1} a \id{O6} \\ \hline
\id{EV-XX} & Identificador de evidencia & \id{EV-01} a \id{EV-13} \\ \hline
\id{ENTR-XX} & Código de participante & \id{ENTR-01} a \id{ENTR-08} \\ \hline
\id{RF-XX} & Requisito funcional & \id{RF-01} a \id{RF-42} \\ \hline
\id{RNF-XX} & Requisito no funcional & \id{RNF-01} a \id{RNF-19} \\ \hline
\id{RD-XX} & Restricción de diseño & \id{RD-01} a \id{RD-10} \\ \hline
\id{RL-XX} & Requisito legal (LOPDP) & \id{RL-01} a \id{RL-08} \\ \hline
\id{CU-XX} & Caso de uso & \id{CU-01} a \id{CU-18} \\ \hline
\id{HU-XX} & Historia de usuario & \id{HU-01} a \id{HU-19} \\ \hline
\id{CA-XX} & Criterio de aceptación & \id{CA-01} a \id{CA-19} \\ \hline
\id{MU-XX} & Prototipo de interfaz & \id{MU-01} a \id{MU-08} \\ \hline
\id{C-XXX} & Componente de la arquitectura & Véase Figura~\ref{fig:componentes} \\ \hline
\id{L-XX} & Limitación declarada & \id{L-01} a \id{L-11} \\ \hline
\id{D-XX} & Desviación del prototipo & \id{D-01} a \id{D-03} \\ \hline
\caption{Diccionario de prefijos de identificadores empleados en el documento}
\end{longtable}

% =====================================================================
\newpage
\section{Repositorio del proyecto}

\subsection{Enlace y estructura}

\begin{center}
\url{https://github.com/gleiston-guerrero/SIMPA_ISR401}
\end{center}

\noindent
Los cinco archivos raíz exigidos ---\id{README.md}, \id{LICENSE}, \id{CITATION.cff},
\id{CHANGELOG.md} y \id{checksums.sha256}--- se encuentran completos en la raíz del repositorio.

\vspace{0.3cm}
\noindent\fbox{\parbox{0.97\textwidth}{%
\small\textbf{Declaración de desviación estructural.} La §8.1 de la guía establece que la raíz del
repositorio debe reproducir el árbol del proyecto. En este repositorio el proyecto reside en la
carpeta \id{AHMRV/}, un nivel por debajo de la raíz.

\vspace{0.2cm}
La desviación es deliberada. La ruta \id{/tree/main/AHMRV} es la declarada en la portada del
documento entregado en el sistema de gestión académica, cuya actividad se encuentra cerrada y no
admite modificación por ningún medio. Trasladar el contenido a la raíz produciría un error 404 en
ese enlace, dado que GitHub no redirige rutas eliminadas, y activaría el gatekeeper G1 con nota
máxima de 2,50 sobre 10.

\vspace{0.2cm}
Ante la disyuntiva entre una desviación estructural penalizable en rúbrica y un enlace roto que
activa un gatekeeper, el equipo optó por preservar la verificabilidad del enlace y declarar la
desviación de forma expresa. La estructura interna de \id{AHMRV/} reproduce íntegramente el árbol
de la §8.1, y el \id{README.md} de la raíz orienta hacia ella.}}

\begin{longtable}{|L{4.6cm}|L{10.8cm}|}
\hline
\rowcolor{grisclaro}\textbf{Carpeta} & \textbf{Contenido} \\
\hline
\endhead
\id{01\_ERS/} & Documento en PDF, fuente \id{.tex}, \id{referencias.bib} e imágenes vectoriales \\ \hline
\id{02\_Evidencias/} & Zona restringida cifrada, consentimientos enmascarados, transcripciones, fotografías, cuestionario, documentos de la organización, actas de validación y codificación temática \\ \hline
\id{03\_Modelado/} & Trece diagramas UML e i* en fuente PlantUML, PNG a 300 dpi y SVG; ocho prototipos de interfaz \\ \hline
\id{04\_Trazabilidad/} & \id{matriz\_trazabilidad.csv} y \id{priorizacion\_moscow\_kano.csv} \\ \hline
\id{05\_MVP/} & Enlace al repositorio del prototipo y video de demostración \\ \hline
\id{06\_Experimento/} & Protocolo, comprobante OSF, instrumentos, consignas de LLM, resultados y scripts \\ \hline
\id{07\_Publicacion/} & Borrador de manuscrito, análisis de revistas y paquete de datos para Zenodo \\ \hline
\id{08\_Etica/} & Paquete ético completo, documentos de Categoría C y adenda de la segunda ronda \\ \hline
\caption{Estructura del repositorio del proyecto}
\end{longtable}

\subsection{Instrucciones de reproducción}

\textbf{Obtención del repositorio con los archivos multimedia:}

\begin{quote}
\id{git clone https://github.com/AlanNVR/}\allowbreak\id{Villafuerte\_Grupo\_AHMRV.git}\\
\id{cd Villafuerte\_Grupo\_AHMRV}\\
\id{git lfs pull}
\end{quote}

\noindent
\textbf{Compilación del presente documento} desde \id{01\_ERS/}:

\begin{quote}
\id{pdflatex ERS\_SRS\_2A\_v1.0.tex}\\
\id{bibtex ERS\_SRS\_2A\_v1.0}\\
\id{pdflatex ERS\_SRS\_2A\_v1.0.tex}\\
\id{pdflatex ERS\_SRS\_2A\_v1.0.tex}
\end{quote}

\noindent
\textbf{Regeneración de los diagramas} desde \id{03\_Modelado/Diagramas\_UML/}:

\begin{quote}
\id{plantuml -tpng -Sdpi=300 -o png *.puml}\\
\id{plantuml -tsvg -o svg *.puml}
\end{quote}

\noindent
\textbf{Verificación de integridad} desde la raíz: \id{sha256sum -c checksums.sha256}

\subsection{Contribución del equipo}

\begin{longtable}{|L{4.4cm}|L{4.0cm}|L{6.8cm}|}
\hline
\rowcolor{grisclaro}\textbf{Integrante} & \textbf{Rol} & \textbf{Contribución principal en la 2A} \\
\hline
\endhead
Villafuerte Rosero Allan Noe & Analista líder & Coordinación, segunda ronda de elicitación, paquete ético y adenda, protocolo experimental \\ \hline
Huilcapi León Denisses Fabiola & Documentadora & Redacción del ERS, migración a \LaTeX, control de versiones \\ \hline
Rizzo Vélez Edson Nagib & Verificador & Verificación de calidad de requisitos, gestión de evidencias y fichas técnicas \\ \hline
Macías Herrera Josthyn Esteban & Modelador & Diagramas UML e i*, matriz de trazabilidad \\ \hline
Arboleda Yanza Francisco Javier & Apoyo modelado & Estructura del repositorio, priorización, diagramas \\ \hline
Alcívar Vélez Anderson Adonis & Apoyo repositorio & Transcripciones, anonimización, prototipo funcional \\ \hline
\caption{Contribución de los integrantes del equipo en la Entrega 4 (2B)}
\end{longtable}

\input{apendice_bdd}

